\section{Geometric depth} \label{XIV.5}

To apply \Ref{XIV.4.2} and its corollaries in practice, one needs a convenient criterion allowing one to verify the \'etale depth hypotheses at the points of $U$. We shall give such a criterion, using the local Lefschetz theorem~\Ref{XIV.4.5}.

\subsection{} \label{XIV.5.1} Let $A$ be a noetherian local ring; when we speak of the \'etale depth of $A$, it will be the depth at the closed point. We shall introduce a notion of ``geometric depth of $A$'', and use \Ref{XIV.4.5} to compare it with the \'etale depth $\profet(A)$.

\begin{proposition} \label{XIV.5.2} Let $A$ be a noetherian local ring; suppose that $A$ is isomorphic to a quotient of a regular local ring $B$ by an ideal $I$ (this is true for example when $A$ is complete, by Cohen's theorem (\textup{$\EGA 0_{\textup{IV}}$~19.8.8})). Let $q$ be the minimal number of generators of $I$; then the number $\dim(B)-q$ is independent of the choice of $B$.
\end{proposition}

The minimal number of generators of $I$ is also equal to the rank of the $k$-vector space $I\otimes_{B}k$, where $k$ denotes the residue field of $A$. One reduces immediately to the case where $A$ is complete, since one has $\widehat{A}\simeq\widehat{B}/\widehat{I}$ with $\dim \widehat{B}=\dim B$ and $\rg_{k}(I\otimes_{B}k)=\rg_{k}(\widehat{I}\otimes_{\widehat{B}}k)$; for the same reason one may suppose that the ring $B$ is complete. Let $B$ and $B^{\prime}$ be two complete regular local rings, $f\colon B\to A$, $f^{\prime}\colon B^{\prime}\to A$ two surjective homomorphisms and $I=\Ker(f)$, $I^{\prime}=\Ker(f^{\prime})$. One must show that
$$
\dim B-\rg_{k}(I\otimes_{B}k)=\dim B^{\prime}-\rg_{k}(I^{\prime}\otimes_{B^{\prime}}k).
$$

Let us first consider the case where one has a factorization of
the form
$$
\xymatrix@C=15pt{ B\ar[rr]^-{f}\ar[rd]_-{g} & & A\\
& B^{\prime}\ar[ru]_-{f^{\prime}} & }$$
with $g$ surjective. Let $J=\Ker(g)$; then $J\subset I$ and $I/J=I^{\prime}$. Since $B^{\prime}$ is regular, $\dim(B^{\prime})=\dim(B)-\rg_{k}(J\otimes_{B}k)$ and $J$ is generated by elements forming part of a regular system of parameters of $B$. It follows that one has the exact sequence
$$0\to J\otimes_{B}k\to I\otimes_{B}k\to J/I\otimes_{B^{\prime}}k\to 0,
$$
and consequently
$$
\dim B-\rg_{k}(I\otimes_{B}k)=\dim B-\rg_{k}(J\otimes_{B}k)-\rg_{k}(J/I\otimes_{B^{\prime}}k)=\dim B^{\prime}-\rg_{k}(I^{\prime}\otimes_{B^{\prime}}k).
$$

The general case reduces to the preceding one; to see this, it suffices to show that one can find a complete regular local ring $B^{\prime\prime}$ and surjective homomorphisms $g\colon B^{\prime\prime}\to B$ and $g^{\prime}\colon B^{\prime\prime}\to B^{\prime}$, making the diagram commutative
\begin{equation*} \label{eq:XIV.5.2.*} \tag{$*$}
\begin{array}{c}
\xymatrix@C=10pt@R=15pt{ & B\ar[rd]^-{f} & \\
B^{\prime\prime}\ar[ru]^-{g}\ar[rd]_-{g^{\prime}} & & A\\
& B^{\prime}.\ar[ru]_-{f^{\prime}} &
}
\end{array}
\end{equation*}
Now, if $W$ is a Cohen ring with residue field $k$, one has a local morphism $W\to A$ which lifts to $B$ and $B^{\prime}$ (\EGA IV~19.8.6), so that one has the commutative diagram
$$
\xymatrix@R=15pt@C=10pt{ & B\ar[rd] & \\
W\ar[ru]\ar[rd] & & A\\
& B^{\prime}.\ar[ru] & }$$
One can find integers $n$ and $n^{\prime}$ and surjective morphisms \hbox{$h\colon W\llbracket T_{1},\dots, T_{n}\rrbracket\!\to\! B$} and $h^{\prime}\colon W\llbracket T_{1}^{\prime},\dots, T_{n^{\prime}}^{\prime}\rrbracket\to B^{\prime}$ which are morphisms of $W$-algebras ($\EGA 0_{\textup{IV}}$ 19.8.8); if one then sets $B^{\prime\prime}=W\llbracket T_{1},\dots, T_{n}, T_{1}^{\prime},\dots, T_{n^{\prime}}^{\prime}\rrbracket$ and if one defines $g$ and~$g^{\prime}$ as morphisms of $W$-algebras such that
$$g(T_{i})=h(T_{i}), \quad g(T_{i}^{\prime})=b_{i}, \quad g^{\prime}(T_{i})=b_{i}^{\prime}, \quad g^{\prime}(T_{i}^{\prime})=h^{\prime}(T_{i}^{\prime}),
$$
where $b_{i}$ (\resp $b_{i}^{\prime}$) is an element of $B$ (\resp of $B^{\prime}$) which lifts $(f^{\prime}\circ h^{\prime})(T_{i}^{\prime})$ (\resp $(f\circ\nobreak h)(T_{i})$), the diagram \eqref{eq:XIV.5.2.*} is indeed commutative.

Proposition \Ref{XIV.5.2} justifies the following definition:

\begin{definition} \label{XIV.5.3}
Let $A$ be a noetherian local ring, $\widehat{A}$ its completion, which is thus isomorphic to the quotient of a complete regular local ring $B$ by an ideal $I$; if $q$ is the minimal number of generators of $I$, one calls the \emph{geometric depth} of $A$ the number
$$
\profgeom(A)=\dim B-q.
$$
\end{definition}

\begin{proposition} \label{XIV.5.4} Let $A$ be a noetherian local ring. Then one has
$$
\profgeom(A)\leq\dim A,
$$
and one has equality if and only if $A$ is a complete intersection.
\end{proposition}

One may suppose $A$ complete. Let then $A=B/I$, where $B$ is a complete regular local ring and $I$ an ideal of $B$. If $(x_{1},\dots, x_{q})$ is a minimal system of generators of $I$, one has $\dim A\geq\dim B-q$, and to say that $\dim A=\dim B-q$ is equivalent to saying that $(x_{1},\dots, x_{q})$ forms part of a system of parameters of $B$ ($\EGA 0_{\textup{IV}}$~16.3.7); the proposition follows immediately.

\begin{proposition} \label{XIV.5.5} Let $A$ and $A^{\prime}$ be noetherian local rings, $f\colon A\to A^{\prime}$ a local homomorphism. Suppose that $f$ is flat and that, denoting by $k$ the residue field of $A$, $A^{\prime}\otimes_{A}k$ is a field, a separable extension of $k$. Then one has
$$
\profgeom(A)=\profgeom(A^{\prime}).
$$
\end{proposition}

Replacing $A$ and $A^{\prime}$ by their completions if necessary, one may suppose $A$ and $A^{\prime}$ complete (it follows from ($\EGA 0_{\textup{III}}$~10.2.1) that the flatness hypothesis is preserved and this is evident for the other hypotheses). Let then $A=B/I$, where $B$ is a regular local ring and $I$ an ideal of $B$. Since $A^{\prime}$ is formally smooth over $A$ ($\EGA 0_{\textup{IV}}$~19.8.2), it follows from ($\EGA 0_{\textup{IV}}$~19.7.2) that one can find a complete noetherian local ring $B^{\prime}$ and a local homomorphism $B\to B^{\prime}$, such that $B^{\prime}$ is a flat $B$-module and that one has $B^{\prime}\otimes_{B}A\simeq A^{\prime}$. One thus has $A^{\prime}\simeq B^{\prime}/IB^{\prime}$; moreover the ring $B^{\prime}$ is regular; indeed, if $\mm$ is the maximal ideal of $B$, $\mm B^{\prime}$ is the maximal ideal of $B^{\prime}$, and, since $\mm$ is generated by a regular sequence by definition of ``regular'', $\mm B^{\prime}$ is generated by a $B^{\prime}$-regular sequence ($\EGA 0_{\textup{IV}}$~15.1.14). Since one evidently has $\dim B=\dim B^{\prime}$, and since $I$ and $IB^{\prime}$ have the same minimal number of generators, the assertion follows.

\refstepcounter{subsection}\label{XIV.5.6}
\begin{enonce*}{Theorem 5.6\ndemark}\ndetext{Illusie has since shown the inequality $\prof_x(\ZZ^/\ell^\nu\ZZ)\geq \profgeom_x(X/S)-\delta(x)+1$ for $x$ a point of $X$ a scheme of finite type over a trait $S$ of residual characteristic prime to $\ell$, and $\nu\geq 1$. If $S$ is of characteristic zero, this is a consequence of theorem~\Ref{XIV.5.6}; see (Illusie~L., ``Perversit\'e et variation'', \emph{Manuscripta Math.} \textbf{112} (2003), p\ptbl 271--295).}
Let $A$ be an excellent local ring of characteristic zero. Then one has
$$
\profet(A)\geq\profgeom(A).
$$
\end{enonce*}

One may suppose $A$ strictly local complete, since geometric depth and \'etale depth are preserved by passage to the strict henselization and to the completion by \Ref{XIV.5.5} and \Ref{XIV.1.16}. Let $A\simeq B/I$, where $B$ is a complete regular local ring, and let $(f_{1},\dots, f_{q})$ be a minimal system of generators of the ideal $I$. One thus has
$$
\pi=\profgeom(A)=\dim B-q.
$$
Consider the closed immersion
$$Y=\Spec A\to X=\Spec B,
$$
and let $U=X-Y=\bigcup_{1\leq i\leq q}X_{f_{i}}$. If $a$ denotes the closed point of $X$, one must show that, for every prime number $p$, one has
$$
\H_{a}^{i}(Y, \ZZ/p\ZZ)=0\quad\textrm{ for }i<\pi.
$$
Since $B$ is regular excellent, one has $\profet(B)=2\dim X$ (\cf \Ref{XIV.1.10}) and consequently $\H_{a}^{i}(X, \ZZ/p\ZZ)=0$ for $i<2\dim X$. It therefore suffices, to prove the theorem, to prove that the morphism
\begin{equation*} \label{eq:XIV.5.6.*} \tag{$*$}
\H_{a}^{i}(X, \ZZ/p\ZZ)\to\H_{a}^{i}(Y, \ZZ/p\ZZ)
\end{equation*}
is bijective for $i<\pi$. For this one applies the local Lefschetz theorem \Ref{XIV.4.5} with $n=\pi+q-1$, $c=q-1$ hence $n-c=\pi$. Note that $U=X-Y$ is the union of the $q$ affine open sets $X_{f_{i}}$. Let us show that one has, for every point $u$ of $U$:
$$
\profet_{u}(X)\geq\pi+q-1-\dim\big(\overline{\{ u\} }\big)=\dim\Oo_{X, u}$$
(where $\overline{\{ u\} }$ denotes the closure of $u$ \emph{in} $X-\overline{\{ a\} }$). Indeed it follows from \Ref{XIV.1.10} that one has
$$
\profet_{u}(X)=2\dim\Oo_{X, u}\geq\dim\Oo_{X, u}.
$$
Using \Ref{XIV.4.5}, one sees that \eqref{eq:XIV.5.6.*} is bijective for $i<\pi$, which completes the proof of the theorem.

\begin{corollary} \label{XIV.5.7} Let $S$ be the spectrum of a field of characteristic zero (\resp a henselian excellent local scheme of characteristic zero), $f\colon X\to S$ a scheme proper over $S$ (\resp over $S-\{ s\}$). Let $U$ be a union of $c+1$ open sets of $X$, \emph{affine}, $Y$ a closed subscheme with underlying space $X-U$, $n$ and $m$ integers $>0$. Suppose that, for every point $u$ of $U$, one has
$$
\profgeom(\Oo_{X, u})\geq n-\dim\big(\overline{\{ u\} }\big)$$
($\overline{\{ u\} }$ the closure of $u$ in $X$). Then the canonical morphism
$$
\H^{i}(X, \ZZ/m\ZZ)\to\H^{i}(Y, \ZZ/m\ZZ)$$
is bijective for $i<n-c-1$, injective for $i=n-c-1$.
\end{corollary}

One applies \Ref{XIV.4.5} and \Ref{XIV.4.6}. The \'etale depth hypotheses at the points of $U$ are satisfied since one has by \Ref{XIV.5.6} $$
\profet_{u}(X)\geq\profgeom(\Oo_{X, u})\geq n-\dim\big(\overline{\{ u\} }\big).
$$
